\chapter{Motivation}

We will begin by studying the solution of the following equation:

\begin{tcolorbox}[colback=gray!5,colframe=gray!80!black]
    \[u_t + uu_x = 0\text{, for }x\in\mathbb{R}\text{ and }t>0\]
\end{tcolorbox}

This is comunly known as the inviscid Burger's equation, a simplified model for fluid mechanics. So, first, we are going to try to find regular solutions for this equation by the method of characteristic.\\
Considering $u(x,t)$ to be a smooth solution for our equation satisfying the initial condition problem

\[u(x,0) = u_0(x)\]

Where $u_0(x)$ is a function of $x$. By defining the chatacteristic curves:

\begin{equation}
    \begin{cases}
        \dot{x}(t) = u(x(t),x)\text{, }x(0) = x_0\\
        \dot{y}(t) = 1\text{, }y(0) = y_0
    \end{cases}
\end{equation}

It is given that $y = t$. By calculating the derivative of $u(x(t),t)$ we obtain: 

\[\dv{u}{t}(x(t),t) = \pdv{u}{t} + \pdv{u}{x}\pdv{x}{t} = 0\]

From this, we can get two conclusions: $u$ is constant along the curve $(x(t),t)$, which implies that $\dot{x}(t) = u_0(x_0)$. By solving this last ODE, we obtain that:

\[x = x_0 + u_0(x_0)t\]

And for all points in this curve $u(x,t) = u_0(x_0)$. By convenience, we will choose this as our initial curve $u_0$:



